\subsection{PRODUCTS AND INVERSE LIMITS OF COMPLETE SPACES}

PROPOSITION 10. Every product of complete uniform spaces is complete. Conversely, if a product of non-empty uniform spaces is complete, then each of the factors is a complete uniform space.

The first assertion is a consequence of the characterization of Cauchy filters and convergent filters on a product space (no. 1, Corollary 2 to Proposition 4 and Chapter I, § 7, no. 6, Corollary 1 to Proposition 10). Conversely, suppose

\[
\mathbf {X} = \prod_ {i \in I} \mathbf {X} _ {i}
\]

is complete (the \(X_{t}\) being non-empty) and let \(F_{x}\) be a Cauchy filter on \(X_{x}\) . For each \(t \neq x\) let \(F_{t}\) be a Cauchy filter on \(X_{t}\) , and consider the product filter (Chapter I, § 6, no. 7)

\[
\mathfrak {F} = \prod_ {i \in I} \mathfrak {F} _ {i} \text { on } X;
\]

\(\mathfrak{F}\) is a Cauchy filter (no. 1, Corollary 2 to Proposition 4), hence is convergent, and therefore so is each \(\mathrm{pr}_{\mathbf{x}}\mathfrak{F} = \mathfrak{F}_{\mathbf{x}}\) (Chapter I, § 7, no. 6, Corollary 1 to Proposition 10).

COROLLARY. Let \((\mathbf{X}_{\alpha}, f_{\alpha\beta})\) be an inverse system of uniform spaces. If the \(X_{\alpha}\) are Hausdorff and complete, then so is \(X = \varprojlim X_{\alpha}\) .

For X is Hausdorff and can be identified with a closed subspace of \(\prod_{i\in I}X_{\alpha}\)

(Chapter I, § 8, no. 2, Corollary 2 to Proposition 7); the corollary therefore follows from Proposition 10 and Proposition 8 (no. 4).

An inverse limit of complete Hausdorff uniform spaces \(X_{\alpha}\) can be empty, even if all the \(X_{\alpha}\) are non-empty and all the \(f_{\alpha\beta}\) are surjective, as is shown by the case of discrete spaces (Set Theory, Chapter III, § 1, Exercise 32). However, we have the following theorem:

THEOREM 1 (Mittag-Leffler). Let \((\mathbf{X}_{\alpha}, f_{\alpha\beta})\) be an inverse system of complete Hausdorff uniform spaces, indexed by a directed set I which has a countable cofinal subset; suppose also that, for each \(\alpha \in I\) , \(X_{\alpha}\) has a countable fundamental system of entourages (*). Finally, suppose that for each \(\alpha \in I\) there is an index \(\beta \geqslant \alpha\) satisfying the following condition:

\((\mathbf{ML}_{\alpha \beta})\) For each \(\gamma \geqslant \beta, f_{\alpha \gamma}(\mathbf{X}_{\gamma})\) is dense in \(f_{\alpha \beta}(\mathbf{X}_{\beta})\) .

Let \(\mathbf{X} = \varprojlim \mathbf{X}_{\alpha}\) and let \(f_{\alpha}\) be the canonical mapping \(\mathbf{X} \to \mathbf{X}_{\alpha}\) . Then for each \(\alpha \in \overleftarrow{\mathbf{I}}\) and each \(\beta \geqslant \alpha\) satisfying \((\mathbf{ML}_{\alpha\beta})\) , \(f_{\alpha}(\mathbf{X})\) is dense in \(f_{\alpha\beta}(\mathbf{X}_{\beta})\) (and consequently \(\mathbf{X}\) is non-empty if the \(\mathbf{X}_{\alpha}\) are all non-empty).

Let \((\lambda_{n})\) be a sequence of indices cofinal in I. Start with an index \(\alpha_0 \in \mathbf{I}\) and define recursively an increasing sequence \((\alpha_{n})\) such that \(\alpha_{n} \geqslant \lambda_{n}\) and such that \((\mathrm{ML}_{\alpha_{n}, \alpha_{n+1}})\) is true. Clearly the sequence \((\alpha_{n})\) is cofinal in I. We shall write \(f_{mn}\) in place of \(f_{\alpha_m\alpha_n}\) for \(m \leqslant n\) , and we shall put \(f_{n,n+1}(\mathbf{X}_{\alpha_{n+1}}) = \mathbf{Y}_n\) . Then, if \(m \leqslant n\) , \(f_{mn}(\mathbf{Y}_n)\) is dense in \(\mathbf{Y}_m\) ; for by definition \(f_{m,n+1}(\mathbf{X}_{\alpha_{n+1}})\) is dense in

\[
f _ {m, m + 1} \left(\mathrm{X} _ {\alpha_ {m + 1}}\right) = \mathrm{Y} _ {m}, \quad \text { and } \quad f _ {m, n + 1} \left(\mathrm{X} _ {\alpha_ {n + 1}}\right) = f _ {m n} \left[ f _ {n, n + 1} \left(\mathrm{X} _ {\alpha_ {n + 1}}\right) \right] = f _ {m n} \left(\mathrm{Y} _ {n}\right).
\]

By induction on \(n\) and \(k\) we can define a fundamental system \((\mathbf{V}_{kn})_{k\in \mathbb{N}}\) of closed symmetric entourages of \(\mathbf{X}_{\alpha_n}\) for each \(n\) such that

\begin{enumerate}
\def\labelenumi{(\arabic{enumi})}
\setcounter{enumi}{1}
\tightlist
\item
\end{enumerate}

\[
\begin{array}{c} \mathbf {V} _ {k + 1, n} ^ {2} \subset \mathbf {V} _ {k n}, \\ (f _ {n, n + 1} \times f _ {n, n + 1}) (\mathbf {V} _ {k, n + 1}) \subset \mathbf {V} _ {k n}. \end{array}\tag{3}
\]

In effect, let \((\mathbf{U}_{kn})_{k\in\mathbb{N}}\) be a fundamental system of entourages of \(X_{\alpha_{n}}\) . If we suppose that the \(V_{kn}\) have been defined for a given n and for all \(k\in N\) , then since \(f_{n,n+1}\) is uniformly continuous we can define the entourage \(V_{k,n+1}\) by induction on k so that (3) is satisfied and

\[
\mathbf {\dot {V}} _ {k + 1, n + 1} ^ {2} \subset \mathbf {V} _ {k, n + 1} \cap \mathbf {U} _ {k + 1, n + 1}
\]

The assertion follows.

Now let \(x_0 \in Y_0\) . We shall show that for each integer \(k > 0\) there is a point \(z \in X\) such that \([x_0, f_{\alpha_0}(z)] \in V_{k-1,0}\) ; this will prove the theorem. Since \(f_{n,n+1}(Y_{n+1})\) is dense in \(Y_n\) , we can define by induction a sequence of points \(x_n \in Y_n\) such that

\[
\left[ x _ {n}, f _ {n, n + 1} \left(x _ {n + 1}\right) \right] \in \mathrm{V} _ {k + n, n}.\tag{4}
\]

By reason of (3) it follows that if \(m \leqslant n\) then

\[
\left[ f _ {m n} (x _ {n}), f _ {m, n + 1} (x _ {n + 1}) \right] \in \mathbf {V} _ {k + n, m}.\tag{5}
\]

From this we conclude that for fixed \(m\) the sequence \((f_{mn}(x_n))_{n\geqslant m}\) is a Cauchy sequence in \(\mathbf{X}_{\alpha_m}\) and therefore converges to a point \(z_m\) ; for by induction it follows from (5) that, for each pair of integers \(p\geqslant m,q > 0\) , we have

\[
\left[ f _ {m p} (x _ {p}), f _ {m, p + q} (x _ {p + q}) \right] \in \mathrm{V} _ {k + p + q - 1, m} \circ \mathrm{V} _ {k + p + q - 2, m} \circ \dots \circ \mathrm{V} _ {k + p, m}\tag{6}
\]

and by virtue of (2) it is clear that the right-hand side of (6) is contained in \(V_{k+p-1,m}\) . Let q increase indefinitely; we infer in particular that, for m=p=0, we have \((x_{0},z_{0})\in\mathbf{V}_{k-1,0}\) , since \(V_{k-1,0}\) is closed. On the other hand, from the relations \(z_{m}=\lim_{n\to\infty}f_{mn}(x_{n})\) and from the continuity of \(f_{m,m+1}\) , we deduce that \(f_{m,m+1}(z_{m+1})=z_{m}\) for each \(m\geqslant0\) . For each \(\gamma\in I\) there is at least one integer n such that \(\alpha_{n}\geqslant\gamma\) ; putting \(z_{\gamma} = f_{\gamma, \alpha_n}(z_n)\) we verify immediately that \(z_{\gamma}\) does not depend on the value of \(n\) such that \(\alpha_n \geqslant \gamma\) , and that the family \((z_\alpha)_{\alpha \in \mathbf{I}}\) so defined is a point \(z\) of \(\mathbf{X} = \varprojlim \mathbf{X}_\alpha\) . Since \(f_{\alpha_0}(z) = z_0\) , the proof is complete.

COROLLARY 1. Let \((\mathbf{X}_{\alpha}, f_{\alpha\beta})\) be an inverse system of sets indexed by a directed set I which has a countable cofinal subset, and suppose that the \(f_{\alpha\beta}\) are surjective. Then if \(X = \varprojlim X_{\alpha}\) , the canonical mapping \(f_{\alpha}: X \to X_{\alpha}\) is surjective for each \(\alpha \in I\) .

Let each \(X_{\alpha}\) carry the discrete uniformity, and apply Theorem 1.

COROLLARY 2. Let I be a directed set which has a countable cofinal subset. Let \((\mathbf{X}_{\alpha}, f_{\alpha \beta})\) and \((\mathbf{X}_{\alpha}^{\prime}, f_{\alpha \beta}^{\prime})\) be two inverse systems of sets indexed by I, and for each \(\alpha \in I\) let \(u_{\alpha}: \mathbf{X}_{\alpha} \to \mathbf{X}_{\alpha}^{\prime}\) be a mapping, such that the \(u_{\alpha}\) form an inverse system of mappings. Let \(u = \varprojlim u_{\alpha}\) . Let \(x = (x_{\alpha}^{\prime})\) be an element of

\[
\mathrm{X} ^ {\prime} = \varprojlim \mathrm{X} _ {\alpha} ^ {\prime}
\]

which satisfies the following condition: for each \(\alpha \in I\) there is an index \(\beta \geqslant \alpha\) such that for all \(\gamma \geqslant \beta\) we have \(f_{\alpha \gamma}(\overline{u}_{\gamma}^{1}(x_{\gamma})) = f_{\alpha \beta}(\overline{u}_{\beta}^{1}(x_{\beta}^{\prime}))\) . Then there is an element \(x \in X\) such that \(u(x) = x'\) .

Apply Theorem 1 to the inverse system of sets \(\overline{u}_{\alpha}^{1}(x_{\alpha}^{\prime})\) , each carrying the discrete uniformity.

* Example. Suppose we are given in C: (i) a sequence \((a_{n})\) of distinct points such that the sequence \((|a_{n}|)\) is increasing and tends to \(+\infty\) ; (ii) for each \(n\) , a rational function \(z \to \mathbb{R}_{n}(z)\) defined in \(\mathbf{C} - \{a_{n}\}\) and having a pole at \(a_{n}\) ; (iii) an increasing sequence \((\mathbf{B}_{n})\) of open discs centred at 0, whose union is \(\mathbf{C}\) , and such that none of the \(a_{k}\) is on the frontier of any of the discs \(\mathbf{B}_{n}\) . For each \(n\) , let \(\mathbf{B}_{n}^{\prime}\) denote the intersection of \(\overline{\mathbf{B}}_{n}\) and the complement in \(\mathbf{C}\) of the set of points \(a_{n}\) ; and let \(\mathbf{X}_{n}\) denote the set of all mappings

\[
z \rightarrow \mathrm{S} (z) = \mathrm{P} (z) + \sum_ {a _ {k} \in \mathrm{B} _ {n}} \mathrm{R} _ {k} (z)
\]

of \(\mathbf{B}_n'\) into \(\mathbf{C}\) , where \(\mathbf{P}\) is the restriction to \(\mathbf{B}_n'\) of a function which is continuous in \(\overline{\mathbf{B}}_n\) and holomorphic in \(\mathbf{B}_n\) . We define a metric in \(\mathbf{X}_n\) by putting

\[
d _ {n} \left(\mathrm{S} _ {1}, \mathrm{S} _ {2}\right) = \sup _ {z \in \mathrm{B} _ {n} ^ {\prime}} | \mathrm{S} _ {1} (z) - \mathrm{S} _ {2} (z) |.
\]

It is easily verified that \(\mathbf{X}_n\) is complete with respect to this metric. Finally, for \(n \leqslant m\) , we define a mapping \(f_{nm}: \mathbf{X}_m \to \mathbf{X}_n\) such that if \(\mathbf{S} \in \mathbf{X}_m\) then \(f_{nm}(\mathbf{S})\) is the restriction of \(\mathbf{S}\) to \(\mathbf{B}_n'\) . It is clear that the \(f_{nm}\) are uniformly continuous and that \((\mathbf{X}_n, f_{nm})\) is an inverse system of uniform spaces. This being so, an element of the inverse limit \(X = \varprojlim X_{n}\) can be canonically identified with a meromorphic function F in C, whose only poles are the points \(a_{n}\) , and which is such that for each n, \(\mathrm{F}(z) - \mathrm{R}_{n}(z)\) is holomorphic at \(a_{n}\) . The classical theorem of Mittag-Leffler asserts that X is not empty; by virtue of Theorem 1, we have only to verify the condition \((\mathrm{ML}_{nn})\) for all n.~Let

\[
\mathrm{S} _ {n} = \mathrm{P} _ {n} + \sum_ {a _ {k} \in \mathrm{B} _ {n}} \mathrm{R} _ {k}
\]

be an element of \(\mathbf{X}_n\) , where \(\mathbf{P}_n\) is continuous in \(\overline{\mathbf{B}}_n\) and holomorphic in \(\mathbf{B}_n\) ; for any \(m \geqslant n\) , let \(\mathbf{Q}_{mn}\) be the restriction of

\[
\sum_ {a _ {h} \in \mathrm{B} _ {m} - \mathrm{B} _ {n}} \mathrm{R} _ {h} \quad \text { to } \quad \mathrm{B} _ {n} ^ {\prime};
\]

this latter sum is a holomorphic function in some neighbourhood of \(\bar{\mathbf{B}}_n\) , hence (by Taylor's theorem) for each \(\varepsilon > 0\) there is a polynomial \(\mathbf{P}_{mn}\) such that \(|\mathbf{Q}_{mn}(\mathbf{Z}) - \mathbf{P}_{mn}(z)| \leqslant \varepsilon\) in \(\mathbf{B}_n\) ; if \(\mathbf{S}_m\) is the restriction of \(\mathbf{S}_n + \mathbf{Q}_{mn} - \mathbf{P}_{mn}\) to \(\mathbf{B}_m'\) , we have \(\mathbf{S}_m \in \mathbf{X}_m\) and \(|\mathbf{S}_m(z) - \mathbf{S}_n(z)| \leqslant \varepsilon\) in \(\mathbf{B}_n'\) . This completes the proof. *

